
\subsubsection{Weight-Space Learning}

The treatment of neural network weights as a learnable data modality has emerged as an active research direction. The 2026 Weight Space Learning Survey categorizes this field into understanding, representation, and generation paradigms. The present work falls under understanding: decoding behavioral properties from weight inspection.

\textbf{Permutation Equivariance.} A fundamental challenge in processing weights is that hidden neurons have no canonical ordering. Zhou et al. introduced Neural Functionals (NF-Layers) that respect this symmetry through parameter-sharing schemes. Their characterization theorem establishes that NF-Layers capture all permutation-equivariant linear maps.

\textbf{Scalable Representations.} SANE addresses scalability by tokenizing weight subsets sequentially. ProbeGen uses learned deep linear probes with shared generators. Set-based Neural Network Encoding handles mixed-architecture model zoos through logit invariance.

\subsubsection{Property Prediction from Weights}

Prior work on property prediction has focused on scalar targets---primarily accuracy and generalization metrics.

\textbf{Weight Statistics.} Unterthiner et al. demonstrated that handcrafted weight statistics achieve $R^2 \approx$ 0.97 for accuracy prediction on the Small CNN Zoo with approximately 30,000 models. The present work tests whether this extends to structured behavioral predictions with a smaller sample of 193 models.

\textbf{Learned Representations.} SANE achieves $R^2$ > 0.9 for accuracy via learned embeddings. Herrmann et al. adapt Deep Weight Space layers for RNNs using functionalist interrogation approaches.

\subsubsection{Behavioral Understanding}

\textbf{Behavioral Loss.} Meynent et al. propose behavioral loss for weight-space autoencoders, showing that structural reconstruction does not guarantee behavioral fidelity. They report 16-20\% accuracy drops despite low reconstruction error, motivating explicit behavioral supervision.
